% Native semantic reconstruction of source physical page 14 (printed page 13).
% Provenance: IEEE Std 1516-2025, continuation of Clause 1.4.3 and Clause 1.4.4.
\HLAStartSemanticPage{14}{13}
\fontsize{10.02pt}{11.523pt}\selectfont
\begingroup
\hbadness=10000\relax
\addfontfeatures{FakeStretch=0.950}
\vspace*{-0.04pt}

expression of its capabilities and data interfaces. The SOM facilitates simulation reuse and provides a\linebreak
starting point for combining the simulation with other participating simulations for a common purpose.\linebreak
Second is a federation object model (FOM). The FOM is a specification defining the information\linebreak
exchanged at runtime to achieve a given set of federation objectives. This information includes object\linebreak
classes, object class attributes, interaction classes, interaction parameters, Management Object Model\linebreak
(MOM) contents, and other relevant information. Third is the MOM, which is a group of predefined\linebreak
constructs that provide support for monitoring and controlling a federation execution. FOMs can be\linebreak
developed from scratch by using relevant elements of the SOMs of the federate applications that will make\linebreak
up the federation or by using existing FOMs or reference FOMs and any combination of these approaches.\linebreak
If existing SOMs or FOMs are used, it is important to understand that there may be parts of these\linebreak
SOMs/FOMs that may not be relevant to other federate applications in the federation and thus will not be\linebreak
incorporated in the FOM. For example, a ship simulation may model boiler pressure, but if a federation\linebreak
needs only to know ship location, boiler pressure may optionally be dropped from the FOM. Also, it may\linebreak
be necessary to reconcile differing units and differing representations among simulations. The OMT\linebreak
provides a standard structure and a set of representational conventions for specifying both SOMs and\linebreak
FOMs.\par
\vspace{18.95pt}

\HLASubclauseHeading[10.02pt]{clause-1.4.4}{1.4.4 Service groups}
\vspace{9.86pt}

The federate interface specification defines the functional interface between federates and the RTI. The RTI\linebreak
provides services to joined federates in a way that is analogous to how a distributed operating system\linebreak
provides services to applications. The services are arranged into seven basic groups as, as follows:\par

\vspace{2.02pt}
\begin{hlaoverviewservicelist}
\item Federation Management: These services allow for the coordination of federation-wide activities\linebreak
throughout the life of a federation execution. Such services include federation execution creation\linebreak
and destruction, federate application joining and resigning, federation synchronization points, and\linebreak
save and restore operations.

\item Declaration Management: These services allow joined federates to specify the types of data they\linebreak
will supply to or receive from the federation execution. This process is done via a set of\linebreak
publication and subscription services along with some related services.

\item Object Management: These services support the life-cycle activities of the objects and interactions\linebreak
used by the joined federates of a federation execution. These services provide for registering and\linebreak
discovering object instances, updating and reflecting the instance attributes associated with these\linebreak
object instances, deleting or removing object instances as well as sending and receiving\linebreak
interactions and other related services. (Note: Formal definitions for each of these terms can be\linebreak
found in the definitions clause of all three HLA specifications.)

\item Ownership Management: These services are used to establish a specific joined federate’s privilege\linebreak
to provide values for an object instance attribute as well as to facilitate dynamic transfer of this\linebreak
privilege (ownership) to other joined federates during a federation execution.

\item Time Management: These services allow joined federates to operate with a logical concept of time\linebreak
and to jointly maintain a distributed virtual clock. These services support discrete event\linebreak
simulations and assurance of causal ordering among events.

\item Data Distribution Management: These services allow joined federates to further specify the\linebreak
distribution conditions (beyond those provided via Declaration Management services) for the\linebreak
specific data they send or ask to receive during a federation execution. The RTI uses this\linebreak
information to route data from producers to consumers in a more tailored manner.

\item Support Services: This group includes miscellaneous services utilized by joined federates for\linebreak
performing such actions as name-to-handle and handle-to-name transformations, the setting of\linebreak
advisory switches, region manipulations, and RTI startup and shutdown.
\end{hlaoverviewservicelist}
\vspace{1.02pt}

The federate interface specification includes programming language-independent definitions of the syntax\linebreak
and semantics of all the services in the interface. It also provides language bindings for Java and C.\par
\endgroup
\clearpage
